\documentclass[11pt]{article}
\usepackage{amsmath, amssymb, amsthm}
\usepackage{geometry}
\usepackage{booktabs}
\usepackage{graphicx}
\usepackage{listings}
\usepackage{xcolor}
\usepackage{hyperref}
\geometry{margin=1.2in}

\lstdefinestyle{cpp}{
  language=C++,
  basicstyle=\ttfamily\small,
  keywordstyle=\bfseries,
  commentstyle=\itshape\color{gray},
  columns=fullflexible,
  keepspaces=true,
  breaklines=true,
  frame=single,
  framerule=0.2pt,
  xleftmargin=1em,
  xrightmargin=1em
}

\title{Movement Metrics Under Grokking\\[0.4em]
\large Instrumenting a Nanda-Style Modular Arithmetic Run\\
with Parameter-Space and Output-Space Trajectory Metrics}
\author{}
\date{}

\begin{document}
\maketitle

\section*{Purpose}

Replicate the canonical grokking setup of Nanda et al.\ (2023) --- modular addition,
one-layer transformer, heavy weight decay, small train fraction --- while continuously
measuring two trajectory metrics: the \emph{parameter-space efficiency ratio} and the
\emph{output-space accord ratio}. The question: do these movement metrics cooccur with,
or precede, the canonical grokking signature (train accuracy saturating early, validation
accuracy jumping late)? A rising accord ratio ahead of the validation breakthrough would
make it a live, loss-independent predictor of circuit formation.

\section{Task and Model}

\paragraph{Task.} $(a + b) \bmod p$ with $p = 113$. The full dataset is all $p^2 = 12{,}769$
ordered pairs. A fixed seeded split takes $3{,}840$ examples ($\approx 30\%$) for training;
the remaining $8{,}929$ are held out for validation. Inputs are the three-token sequence
$[a,\, b,\, =]$, one-hot over a vocabulary of size $p + 1$; the target is the one-hot answer
read at the $=$ position.

\paragraph{Model.} One pre-norm transformer block:
token embedding ($\mathbb{R}^{114} \to \mathbb{R}^{128}$ per position), learned positional
embedding, multi-head self-attention (4 heads, head dim 32), MLP ($128 \to 512 \to 128$, ReLU),
position readout at $=$, linear unembedding to $p$ logits. Softmax cross-entropy loss.

\paragraph{Optimizer.} AdamW: $\beta_1 = 0.9$, $\beta_2 = 0.98$, learning rate $10^{-3}$,
decoupled weight decay $\lambda_{\mathrm{wd}} = 1.0$. Weight decay is essential: grokking
largely disappears without it. Minibatches of 128; $5 \times 10^4$ steps.

\section{Expression in TTTN}

The experiment is implemented in TTTN, a header-only C++20 library in which tensor
\emph{shapes are types}: a rank-$r$ tensor of dimensions $d_1 \times \cdots \times d_r$
is the type \texttt{Tensor<$d_1$,\dots,$d_r$>}, and every contraction, layer, and network
carries its full shape signature at compile time. Three consequences matter here.

\paragraph{Networks are types.} A layer is any type satisfying the \texttt{Block} concept:
it declares \texttt{InputTensor} and \texttt{OutputTensor}, provides shape-typed
\texttt{Forward}/\texttt{Backward}, and exposes its parameters through \texttt{all\_params()}.
A network is a \texttt{TrainableTensorNetwork<Blocks...>} --- a compile-time-checked chain in
which each block's \texttt{OutputTensor} must equal the next block's \texttt{InputTensor}
(enforced by \texttt{static\_assert}; an ill-shaped architecture does not compile).
The model of Section~1 is, verbatim from the runner:

\begin{lstlisting}[style=cpp]
using EmbedBlock   = MapDenseMDBlock<Tensor<3,114>, Tensor<128>, 1>;
using PosEmbBlock  = PositionalEmbeddingBlock<3, 128>;
using TxfBlock     = TransformerBlock<Tensor<3,128>, 4, 512>;
using SelectBlock  = SelectPositionBlock<3, 128, 2>;
using UnembedBlock = DenseMDBlock<Tensor<128>, Tensor<113>>;

using NandaNet = TrainableTensorNetwork<
    EmbedBlock, PosEmbBlock, TxfBlock, SelectBlock, UnembedBlock>;
\end{lstlisting}

The gradient-descent program one is really writing must type-check in the language of
linear algebra; here it type-checks in C++ as well, and the two checks are the same check.

\paragraph{The parameter walk is free.} Because every block exposes \texttt{all\_params()}
as a tuple of \texttt{Param} objects and tuples concatenate structurally, the entire network
flattens to one canonical parameter order at compile time
(\texttt{NandaNet::TotalParamCount} is a \texttt{constexpr}). Every metric in this document
--- the per-parameter accumulators $G_i, D_i$, the flat $\Delta\theta$ snapshot, the Jacobian
column extraction, the output-space accumulators indexed by $i$ --- is a single fold over
this tuple. No name matching, no reflection, no bookkeeping: the type system \emph{is} the
bookkeeping.

\paragraph{Instrumentation lives where the physics happens.} Each \texttt{Param} owns its
value, gradient, Adam moments, and trajectory accumulators side by side; $G_i$ and $D_i$ are
updated inside the optimizer step itself, so parameter-space movement is captured at the exact
moment it occurs, at every step, for free. The Jacobian computation reuses the ordinary
backward infrastructure --- a unit vector $e_j$ passed as the upstream gradient --- so the
output-space metrics require no second autodiff system: the same typed backward pass that
trains the network also measures it.

\section{The Base Accumulation: Parameter-Space Movement}

Let $\theta_i(t)$ be parameter element $i$ after step $t$, and
$\Delta\theta_i(t) = \theta_i(t) - \theta_i(t-1)$ the realized update (Adam step
plus weight-decay shrink). Two running accumulators per parameter:
\begin{align}
    G_i &= \sum_t \bigl|\Delta\theta_i(t)\bigr|
        &&\text{(gross path --- the odometer)} \\[2pt]
    D_i &= \sum_t \Delta\theta_i(t) \;=\; \theta_i(T) - \theta_i(0)
        &&\text{(net displacement)}
\end{align}

\paragraph{Efficiency ratio.}
\begin{equation}
    \eta_{\mathrm{param}} \;=\; \frac{\|D\|_2}{\sum_i G_i} \;\in\; [0, 1]
\end{equation}
$\eta_{\mathrm{param}} = 1$ iff every step of every parameter moves monotonically toward its
final position (geodesic motion); $\eta_{\mathrm{param}} \to 0$ under pure churn. This metric
treats all parameters as equally important and all directions in parameter space as
equivalent --- it is blind to what the movement \emph{does}. The output-space metrics below
correct for this.

\section{Output-Space Movement: The Jacobian as the Metric}

\paragraph{The Jacobian.} For the network function $F : \mathbb{R}^n \to \mathbb{R}^p$
(logit vector, $p = 113$) evaluated at a fixed reference input $x$,
\begin{equation}
    J(t) \;=\; \frac{\partial F(x)}{\partial \theta}\bigg|_{\theta(t)} \;\in\; \mathbb{R}^{p \times P},
    \qquad
    J_i(t) \;=\; \frac{\partial F(x)}{\partial \theta_i}\bigg|_{\theta(t)} \;\in\; \mathbb{R}^{p}
\end{equation}
Column $J_i$ is parameter $i$'s \emph{complete first-order influence} on the output:
perturbing $\theta_i$ by $\varepsilon$ moves the output by $\varepsilon\, J_i + O(\varepsilon^2)$.
Every parameter's influence is expressed in the same coordinates --- the model's output shape.
$J$ is computed by $p$ reverse-mode passes (unit vector $e_j$ as upstream gradient collects row
$j$ for all parameters at once), averaged over a small fixed batch of reference inputs, and is
recomputed at every training step.

\paragraph{Output-space GROSS and NET.} Each step contributes its realized movement,
\emph{metricized by the Jacobian at that moment}, to two running output-shape accumulators:
\begin{align}
    \mathrm{NET}_j &= \sum_t \sum_i \Delta\theta_i(t)\, J_{ij}(t)
        \;=\; \sum_t \bigl(J(t)\,\Delta\theta(t)\bigr)_j \\[2pt]
    \mathrm{GROSS}_j &= \sum_t \sum_i \bigl|\Delta\theta_i(t)\, J_{ij}(t)\bigr|
\end{align}
$\mathrm{NET} \in \mathbb{R}^p$ is the accumulated first-order output displacement --- the signed
sum of what every update did to the function. $\mathrm{GROSS} \in \mathbb{R}^p$ is the same sum
with all cancellation unfolded: the total output-space work performed, counting fights between
parameters (and reversals across time) at full price.

Per-parameter versions of both accumulators
($\mathrm{net}_i,\, \mathrm{gross}_i \in \mathbb{R}^p$ for each $i$) are also maintained and
serialized at checkpoints, supporting post-hoc leverage-weighted analysis.

\paragraph{Accord ratio.}
\begin{equation}
    \eta_{\mathrm{out}} \;=\; \frac{\|\mathrm{NET}\|_2}{\|\mathrm{GROSS}\|_2} \;\in\; [0, 1]
\end{equation}
$\eta_{\mathrm{out}} = 1$ iff every per-step, per-parameter contribution
$\Delta\theta_i(t)\,J_i(t)$ pointed in a consistent direction per output coordinate ---
perfectly collaborative movement. Departure from 1 measures destructive interference:
parameters (or time steps) doing output-space work that cancels.

\paragraph{Per-dimension accord.} The ratio disaggregates over output coordinates:
\begin{equation}
    \eta_j \;=\; \frac{|\mathrm{NET}_j|}{\mathrm{GROSS}_j}
\end{equation}
revealing \emph{which logits} receive collaborative updates and which receive churn.
Under the known Fourier mechanism of modular addition, a natural refinement is to ask whether
accord rises first in the output directions associated with the network's key frequencies.

\section{Relation to the Fisher Information Matrix}

The empirical Fisher at $x$ is $\mathcal{F} = J^\top J \in \mathbb{R}^{P \times P}$, with
\begin{equation}
    \mathcal{F}_{ii} = \|J_i\|_2^2, \qquad
    \mathcal{F}_{ik} = \langle J_i, J_k \rangle .
\end{equation}
The diagonal treats parameters as if they influenced the output orthogonally; the off-diagonal
encodes how parameters share output directions --- the collaboration structure. The identity
\begin{equation}
    \|J\,\Delta\theta\|_2^2
    \;=\; \Delta\theta^\top \mathcal{F}\, \Delta\theta
    \;=\; \underbrace{\sum_i \Delta\theta_i^2\, \mathcal{F}_{ii}}_{\text{diagonal}}
    \;+\; \underbrace{2 \sum_{i<k} \Delta\theta_i\, \Delta\theta_k\, \mathcal{F}_{ik}}_{\text{off-diagonal (collaboration)}}
\end{equation}
shows that the accord ratio probes the full off-diagonal structure of $\mathcal{F}$ under the
realized update directions --- without ever materializing the $P \times P$ matrix
($\sim 10^{10}$ entries at $P \sim 10^5$). Only $J$ itself ($p \times P$) is formed, used once
per step, and discarded.

\section{Measurement Protocol}

Per training step:
\begin{enumerate}
    \item Snapshot $\theta$ (flat, all parameters).
    \item One AdamW training step on a minibatch; recover $\Delta\theta(t)$ by subtraction.
          (Parameter-space accumulators $G_i, D_i$ update inside the optimizer step itself.)
    \item Compute $J(t)$ at the current weights on the fixed reference batch
          ($p \times |\mathcal{B}_{\mathrm{ref}}|$ backward passes).
    \item Accumulate $\Delta\theta_i(t)\,J_i(t)$ into the per-parameter and network-level
          output-space GROSS/NET tensors.
\end{enumerate}
Every 25 steps: evaluate loss and accuracy on the full train and validation splits; log all
scalar metrics to CSV. Every 500 steps: serialize weights, optimizer state, trajectory
accumulators, and the full per-parameter output-space tensors (resumable, crash-safe).

\section{What to Look For}

\begin{center}
\renewcommand{\arraystretch}{1.4}
\begin{tabular}{lp{8.5cm}}
\toprule
\textbf{Observation} & \textbf{Interpretation} \\
\midrule
$\eta_{\mathrm{out}}$ rises \emph{before} val-acc jump &
    Accord is a leading indicator of circuit formation; collaborative output-space movement
    precedes behavioral generalization. \\
$\eta_{\mathrm{out}}$ rises \emph{with} val-acc jump &
    Accord cooccurs with grokking; still a useful loss-independent correlate. \\
$\eta_{\mathrm{out}}$ flat through grokking &
    Clean negative result: accord at the network level does not resolve the transition;
    per-dimension or per-parameter disaggregation required. \\
$\eta_{\mathrm{param}}$ vs.\ $\eta_{\mathrm{out}}$ divergence &
    Movement directed in parameter space but conflicted in function space (or vice versa) ---
    the two geometries decouple, and when they do is itself informative. \\
$\eta_j$ rises first at key-frequency logits &
    Accord localizes to the known Fourier circuit --- strongest possible link between the
    movement metrics and the mechanistic account. \\
\bottomrule
\end{tabular}
\end{center}

\section{Results}

% ── dashboard figure ──
% \begin{figure}[h]
%     \centering
%     \includegraphics[width=\textwidth]{checkpoints_nanda_grokking/dashboard.png}
%     \caption{Grokking signature (top), parameter-space efficiency and output-space accord
%     (middle), cumulative output-space movement (bottom). Dashed line: first step with
%     validation accuracy $\geq 90\%$.}
% \end{figure}

% TODO: results prose.

\section*{Implementation}

TTTN (header-only C++20). Metrics: \texttt{src/NetworkUtil.hpp} (per-parameter $G_i, D_i$
accumulated inside \texttt{Param::update}), \texttt{src/FunctionalInfluence.hpp}
(\texttt{ComputeJacobian}, \texttt{OutputSpaceTrajectory}). Runner:
\texttt{nanda\_grokking.cpp}. Dashboard: \texttt{tools/nanda\_grokking\_dashboard.py}.
Framework: \texttt{PARAMETER\_MOVEMENT.md}.

\begin{thebibliography}{9}
\bibitem{nanda2023}
Nanda, N., Chan, L., Lieberum, T., Smith, J., Steinhardt, J.
\emph{Progress measures for grokking via mechanistic interpretability.} ICLR 2023.

\bibitem{power2022}
Power, A., Burda, Y., Edwards, H., Babuschkin, I., Misra, V.
\emph{Grokking: Generalization beyond overfitting on small algorithmic datasets.} 2022.
\end{thebibliography}

\end{document}
